\documentclass[runningheads]{llncs}

\usepackage[T1]{fontenc}
\usepackage[utf8]{inputenc}
\usepackage{lmodern}
\usepackage{amsmath}
\usepackage{graphicx}
\usepackage{array}
\usepackage{booktabs}
\usepackage{xcolor}
\usepackage{tikz}
\usetikzlibrary{arrows.meta,positioning,fit,backgrounds,calc}
\usepackage{url}
\usepackage{hyperref}
\renewcommand\UrlFont{\color{blue}\rmfamily}
% This llncs.cls copy does not define \orcidID; provide a lightweight version.
\providecommand{\orcidID}[1]{\href{https://orcid.org/#1}{\textsuperscript{[#1]}}}

\begin{document}

\title{Work-in-Progress: From Generic Chatbots to Teacher-Configurable AI Tutors: The Design and Evaluation of Novedu in Secondary Education}
\titlerunning{Work-in-Progress: Teacher-Configurable AI Tutors with Novedu}

\author{Rainer Stropek\orcidID{0009-0003-0970-1663} \and Andreas Probst\orcidID{0000-0002-6699-0866} \and Peter Bauer\orcidID{0009-0001-7200-7143} \and Lukas Larndorfer \and Melanie Bauer}
\authorrunning{R. Stropek et al.}

\institute{HTL Leonding, Leonding, Austria\\
\email{r.stropek@htl-leonding.ac.at}}

\maketitle

\begin{abstract}
Secondary-school students increasingly use general-purpose chatbots for learning, but these systems are rarely aligned with classroom context, teacher intentions, or a school's privacy requirements. This work-in-progress paper presents \emph{Novedu}, a teacher-configurable AI-tutoring platform developed at HTL Leonding, an Austrian technical upper-secondary school. Our position is that a school cannot realistically align large language model (LLM) tutors by training models, but can align them institutionally through teacher-owned configuration, curricular grounding, privacy defaults, and evaluation: tutors are assembled from reusable, teacher-authored modules and served through in-region, open-weight models. A requirements survey ($n=318$) yields four central requirements (curricular grounding, factual reliability, configurable scaffolding, and aggregated oversight), each traced to a design response; a reproducible evaluation harness gives preliminary one-topic evidence that an in-region open-weight model need not fall behind frontier baselines, while depending strongly on teacher-authored prompt detail. The platform and classroom evaluation remain ongoing.
\keywords{generative AI in education $\cdot$ teacher-configurable AI tutors $\cdot$ pedagogical alignment $\cdot$ in-region open-weight LLMs}
\end{abstract}

\section{Introduction}
\label{sec:intro}

Generative AI has become an everyday tool in secondary education. At our school adoption is near-universal: $94\%$ of students have used an AI tool, $84.5\%$ regularly (Chapter~\ref{sec:requirements}). The opportunities and risks are well rehearsed~\cite{kasneci2023chatgpt}; what remains unresolved for an individual school is not \emph{whether} students use AI, but \emph{how} the institution can channel that use toward learning while meeting its legal and pedagogical obligations.

For an institution, a general-purpose chatbot is a poor default for two distinct reasons: one pedagogical, one institutional. Pedagogically, used as-is a large language model (LLM) tends to provide immediate answers rather than guiding students through problem solving~\cite{sonkar2024pedagogical}, and it knows neither the class's current topic nor the teacher's intentions. Institutionally, sending student data to commercial cloud LLMs raises data-protection obligations under the General Data Protection Regulation (GDPR) and, for some educational uses, governance obligations under the EU AI Act, and it gives teachers no oversight of its use in their subject.

This paper presents \emph{Novedu}, a teacher-configurable AI-tutoring software developed as a cooperative student software-engineering project at HTL Leonding, a school of higher technical vocational education in Austria with a strong informatics profile, and deployed at the school. Our central position is that a school cannot realistically align tutors by \emph{training} models, but it can align them \emph{institutionally}: by giving teachers ownership of configuration, curricular grounding, deployment, privacy defaults, and evaluation.

The paper is organized around three research questions, which map onto its three contributions: \textbf{RQ1}: What do students and teachers at a technically oriented secondary school require of an institutional AI tutor? (a requirements evidence base traced to design decisions); \textbf{RQ2}: How can a school align tutors with those requirements through teacher-owned configuration rather than model training? (a teacher-configurable architecture over in-region, open-weight models); \textbf{RQ3}: Can a small, in-region open-weight model deliver the intended Socratic tutoring behavior, and how much prompt detail must teachers provide? (a reproducible evaluation path with preliminary results).

\section{Related Work and Positioning}
\label{sec:related}

Recent work converges on one idea, \emph{pedagogical alignment}: a general-purpose LLM, used as-is, behaves as a generic information provider that hands over answers, whereas effective teaching requires scaffolded, Socratic guidance~\cite{sonkar2024pedagogical}; the works differ in \emph{how} they close that gap. One family aligns models through training: Sonkar et al.~\cite{sonkar2024pedagogical} use preference learning to make a model guide rather than answer, and contribute metrics that quantify scaffolding versus direct answers. A second family aligns through prompt or configuration: Gabrov\v{s}ek and Rihtar\v{s}i\v{c}~\cite{gabrovsek2025custom} deploy a configurable generative-AI tutor in an undergraduate STEM lab and report that structured, declarative prompt styles were rated highest by students.

Novedu adopts the configuration route, the most realistic lever available to a secondary school, which cannot curate preference datasets or fine-tune models but can author prompts and use in-region open-weight models. A generic LLM with careful prompt engineering can of course be made context-aware; the open question is \emph{who} does that engineering and how it is shared, versioned, and checked. Where Gabrov\v{s}ek and Rihtar\v{s}i\v{c} tune one tutor for one course, we study \emph{the school's path to fielding tutors}: an institutional embedding of teacher ownership, curricular fit, privacy, in-region inference, and evaluation. The research questions follow from this positioning: RQ1 asks what such an embedding must satisfy, where the literature offers student ratings of tutors but little requirements evidence from secondary schools; RQ2 asks how teachers, not developers, can tune existing LLMs toward their learning objectives; RQ3 asks whether the result holds the Socratic stance on models a school can run, building on work measuring the \emph{pedagogical} quality of responses rather than answer correctness~\cite{tack2022aiteacher} and on the LLM-as-a-judge paradigm and its biases~\cite{zheng2023judging}.

\section{Requirements and Design Process}
\label{sec:requirements}

\paragraph{From exploratory use to a scope-reduced platform.}
Novedu emerged from a requirements-driven progression: a vision document committed to teacher-edited prompts, per-tutor model choice, open-source components, and self-hostable hosting. Students then composed open-source products (Open WebUI with a model gateway, school single sign-on, per-user budgets) into a working environment (Fig.~\ref{fig:approaches}a); this first version leaned on retrieval-augmented generation (RAG)~\cite{lewis2020rag} to bring curriculum material into the tutor; rather than being evaluated as a product, its use informed the requirements process and motivated the scope reduction reported here. Guided by a senior product owner, the team produced a requirements specification and a user-experience (UX) prototype of an institutional platform; a cross-disciplinary teacher panel received it well but judged it too large to deploy quickly, so the vision was reduced to a minimal deployable tutor workflow, deferring admin consoles and user/group management. Students built the prototype, the requirements and UX artifacts, and parts of the continuation; the minimum viable product (MVP) seed was created in a teacher-led workshop.

\paragraph{Requirements survey.}
Alongside this process, a two-variant questionnaire (students and teachers), fielded in supervised in-class sessions as a pretest, yielded $318$ valid responses ($284$ students, $34$ teachers) on a 4-point Likert scale with no neutral midpoint (theoretical mean $2.5$). The most important caveat is sample bias: the school is informatics-heavy ($72.9\%$ of student respondents in the IT branch), so the findings do not generalize directly to general-education schools. Four requirements dominate across both groups. \emph{Curricular grounding} is the strongest single signal: the top student item is that the tutor know their concrete lesson content ($M=3.53$). \emph{Factual reliability} is critical on both sides: wrong answers are students' most frequent objection and teachers' top property ($84.8\%$). \emph{Configurable scaffolding} beats any fixed mode: students' top explanation choice is ``it depends, let me choose'' ($52\%$). \emph{Non-surveillance oversight} is required by both: students dislike teachers reading even anonymized chats ($57\%$), while teachers want aggregated statistics ($75.8\%$) over individual chats ($6\%$). Two adoption findings frame the rest: teachers are willing but under-prepared (only $6\%$ feel well-informed on the legal framework; training demand $M=3.27$), and willingness to author tutors is conditional on low effort.

Table~\ref{tab:trace} traces each requirement to a design response (Chapter~\ref{sec:design}); Chapter~\ref{sec:limitations} discusses the gaps that remain.

\begin{table}[t]
\centering
\caption{Requirements-to-design traceability.}
\label{tab:trace}
\renewcommand{\arraystretch}{1.0}
\footnotesize
\begin{tabular}{>{\raggedright\arraybackslash}p{3.6cm}>{\raggedright\arraybackslash}p{3.4cm}>{\raggedright\arraybackslash}p{3.9cm}}
\toprule
\textbf{Survey signal} & \textbf{Design implication} & \textbf{Novedu response} \\
\midrule
Tutor must know lesson content ($M{=}3.53$) & Curricular grounding $>$ generic chat & Teacher-authored topic / tutor modules \\
Wrong answers top objection & Reliability must be evaluated & SocraticTutorEval harness (Chap.~\ref{sec:eval}) \\
Want configurable help style ($52\%$) & One Socratic mode is too rigid & Reusable YAML behavior fragments \\
Reject read-along; want aggregate stats & Oversight must be non-surveillance & Anonymous defaults, aggregate code stats \\
Want training and templates & Adoption depends on low effort & Prompt library, starter tutors, AI authoring aids \\
\bottomrule
\end{tabular}
\end{table}

\section{Novedu Design}
\label{sec:design}

\subsection{Teacher-Configurable Tutors as Modular Configuration}
\label{sec:tutormodel}

A single idea organizes the architecture: a tutor is a plain, human-readable file, and the application is a configuration-driven layer around it, so pedagogy lives in teacher-authored configuration, not in application code or model weights. Two kinds of YAML files are composed. \emph{Fragment libraries} are centrally maintained, reusable prompt modules (a persona, ground rules such as a ``never give away the solution'' Socratic policy, a topic description), each versioned and parameterized by an input schema. \emph{Tutor files}, written per subject, select fragments, fill in their inputs, add final instructions, and reference a model. A complete tutor file thus fits on a dozen lines naming a model, referencing fragments (for example a Socratic no-solution policy and a topic), and adding final instructions; full examples are in the project repository. At request time a framework-agnostic pipeline fetches, validates, and assembles the fragments into a system prompt for a single agent answering over an in-region open-weight model. Unlike the prototype, the MVP carries no retrieval pipeline: curriculum context lives in teacher-authored fragments, trading automatic recall for teacher control and simpler deployment. This pedagogical content is produced by a cross-disciplinary, teacher-owned effort: a shared, version-controlled library of tutor fragments authored by teachers of German, English, mathematics, biomedicine, and programming, so the same behavior module serves all subjects alike.

\paragraph{Authoring usability.}
A first presentation to non-technical teachers showed that hand-authoring YAML is a real barrier, matching the survey's low-effort precondition. We address this without abandoning the plain-file model. The intended teacher workflow needs no YAML expertise or developer tooling: a teacher starts from a starter tutor or the prompt library, describes changes in prose to an AI authoring aid that emits a valid tutor definition, and saves it inside the app, where it is validated and versioned (no Git or terminal involved); the teacher then mints a code, tries the tutor as a student would, and iterates, maintaining it later by the same route. Technically inclined teachers can also use a JSON Schema for editor autocompletion and a command-line checker. A graphical editor that hides YAML entirely is planned (a final-year student project).

\subsection{One Abstraction, Several Activities}
\label{sec:activities}

The configurable-file idea generalizes beyond chat: a teacher mints a short \emph{code} pointing at an activity definition, and one access-gating and statistics mechanism serves whichever activity the code names. Fig.~\ref{fig:approaches} contrasts this with the prototype, which composed a generic chat front end (Open WebUI) with a separate LLM proxy (LiteLLM); Novedu integrates both into one teacher-configurable app serving all four activity types. Tutors carry the most pedagogical weight; quizzes pair open-ended questions with a server-only grading rubric scored by a stateless grader; writing provides a split-screen coach that can read but not edit the student's draft; and coding injects a teacher-authored system prompt into an OpenAI-compatible proxy for students' coding agents. Anonymity is the default except where graded writing requires attribution. Status: tutors, quizzes, writing, in-app authoring with validation, codes, and aggregate statistics are operational and in use at the school; the coding proxy is implemented but not yet in production; the graphical editor, curriculum retrieval, admin consoles, and smartphone support are planned.

\begin{figure}[t]
\centering
\includegraphics[width=0.6\textwidth]{figures/NoveduApproachesComparison}
\caption{From composed products to an integrated platform: (a) the exploratory prototype; (b) the Novedu MVP serving the four activity types (Chapter~\ref{sec:activities}).}
\label{fig:approaches}
\end{figure}

\subsection{Self-Hosting and Privacy Defaults}
\label{sec:privacy}

Institutional fit rests on a few deployment decisions. Inference avoids commercial cloud LLMs: open-weight, quantized models (roughly 27--31 billion parameters) are served behind an OpenAI-compatible API, addressed server-side only, by an in-region GPU server operated by the Software Competence Center Hagenberg (SCCH), an Austrian research center, through its open-source \emph{LLM:2go} offering, provided to the school at no cost within a government-funded research project. Identity is the school's single sign-on (teacher actions gated by group membership), and privacy is the default: activities are anonymous unless a teacher opts out, the quiz grader stores nothing, telemetry carries no message content, and chat history is not retained. We claim privacy-\emph{aware} design, not demonstrated legal compliance: these choices reduce exposure by avoiding data-collecting LLM services and keeping processing in-region, but SCCH is a separate processor, and a data-processing agreement, an AI-Act risk classification, and a full compliance assessment are still pending. We also note that institutional licences of commercial assistants under enterprise data-protection terms (for example Microsoft 365 Copilot~\cite{microsoft2025edp}) can offer contractual guarantees of their own; Novedu's in-region choice rests as much on per-tutor model choice, evaluability, and cost as on privacy.

\section{Preliminary Evaluation}
\label{sec:eval}

The survey measures stated preferences. To measure \emph{behavior}, whether a configured tutor actually withholds the answer, we built \emph{SocraticTutorEval}: given a partial student--tutor conversation ending on a student turn, it generates the next tutor reply with several LLMs and has an LLM judge score each against a per-scenario rubric. A rubric states the scenario's goal, positive criteria (e.g., ``asks one guiding question about the student's code'') and negative criteria (e.g., ``states the corrected code''). The judge prompt presents system prompt, conversation, reply, and rubric and requests structured output: per criterion a met/violated verdict with a quoted evidence span, a goal-alignment score, and any Socratic violations from a fixed list (gave answer, multiple questions, lectured, judgmental, off-topic). A pass requires no negative-criterion violation and $\ge 70\%$ of positive criteria met; disagreement between this rule and the judge's own verdict is flagged for rubric review. A failing reply does not indicate a wrong answer but a departure from the behavior the teacher specified. We use it to probe RQ3 on one programming topic (linked lists): three system prompts of increasing length for the same tutor (a 6-line \emph{minimal} prompt of two sentences, an 84-line \emph{short}, and a 172-line \emph{full} prompt), across one in-region open-weight model (the SCCH endpoint) and four frontier baselines (Table~\ref{tab:harness}). Each cell aggregates $36$ judged conversations (12 scenarios $\times$ 3 replies), each scored once by \texttt{gpt-5.4} at high reasoning effort.

\begin{table}[t]
\centering
\caption{Pass rate (\%) and raw passes out of 36 judged replies by model and prompt length (linked-lists scenarios). With $n=36$ per cell, 95\% Wilson intervals span about $\pm16$ points, so within-column differences are not statistically separable.}
\label{tab:harness}
\begin{tabular}{lccc}
\toprule
\textbf{Model} & \textbf{6-line} & \textbf{84-line} & \textbf{172-line} \\
\midrule
gemma-4-31B-FP8 (SCCH, in-region) & \textbf{47.2} (17) & 55.6 (20) & \textbf{61.1} (22) \\
gpt-oss-120b (open weight, medium) & 33.3 (12) & 52.8 (19) & 61.1 (22) \\
gpt-5.4-mini & 27.8 (10) & 55.6 (20) & 61.1 (22) \\
gpt-5.4 (medium reasoning) & 33.3 (12) & 47.2 (17) & 55.6 (20) \\
gpt-5.4 (high reasoning) & 27.8 (10) & 50.0 (18) & 47.2 (17) \\
\bottomrule
\end{tabular}
\end{table}

Two findings stand out. First, \emph{prompt detail matters substantially} (RQ2): moving from the full prompt to the two-sentence minimal prompt lowers the pass rate for every model, by $14$ points for the in-region model and up to $33$ points for \texttt{gpt-5.4-mini}, reinforcing the need for authoring support (Chapter~\ref{sec:tutormodel}). Second, on this Socratic-adherence metric a small \emph{in-region open-weight} model is viable (RQ3): \texttt{gemma-4-31B} reaches the same pass count as the best frontier configurations on the full prompt ($22/36$) and is descriptively the most robust to a terse prompt ($47.2\%$, versus $\le 33.3\%$); higher reasoning effort did not help. Given the interval widths, this shows the in-region model did not fall behind in this sample, not that the models are equivalent. The frontier models do not score lower because they are weaker: their characteristic failure is over-responding, handing the student a full answer and breaking the Socratic constraint. These preliminary results make in-region open-weight inference plausible for further classroom testing without establishing equivalence with, let alone superiority over, commercial models.

These results are preliminary: a single topic, small per-cell counts, and an LLM judge whose biases and calibration require scrutiny~\cite{zheng2023judging} and may diverge from human teachers on helpfulness~\cite{tack2022aiteacher}. The harness has been shared with the teacher group and first cross-teacher experiments have begun; validating the judge against independent teacher ratings is part of the ongoing, multi-subject analysis. Model IDs, decoding parameters, judge prompt, and tutor prompts are archived with the harness.

\section{Limitations and Next Steps}
\label{sec:limitations}

The platform is a learning prototype: functional rather than polished, not yet usable on smartphones, and single-tenant. The evaluation is correspondingly bounded: the survey is a single-site, supervised, self-report pretest, with an informatics-heavy sample that limits external validity; the harness covers one topic and depends on an LLM judge that needs calibration; and we have no classroom learning-outcome data yet. Because students reject read-along on individual chats even when anonymized ($57\%$), individual-level log analysis as in~\cite{gabrovsek2025custom} is hard to justify; log-based work must use aggregated data. We will broaden the questionnaire to additional schools in the coming academic year to remove the branch bias of the pretest.

The most important gap is that our strongest requirement outruns the implementation: students demand that the tutor know their concrete lesson material, yet the MVP supplies this only through teacher-authored prompt modules, not retrieval over the school's learning-management system. Closing that gap is the priority next step, but not necessarily via RAG: we plan to test whether teacher-authored prompts, optionally extended with a lightweight tool-based (grep-style) lookup over curriculum material, are competitive with a full RAG pipeline, since recent evidence finds no universal winner between retrieval and long-context prompting~\cite{li2025lara} and documents characteristic RAG failure modes~\cite{agrawal2024mindfulrag}. Adoption is equally central: the graphical editor and teacher training are conditions for success. Further work includes validating the LLM judge against independent teacher ratings of judged replies across subjects, a behavioral, in-classroom evaluation that compares classes using a configured tutor with classes using a generic chatbot on aggregated, anonymized data, and shipping the coding proxy; the cooperative-development model and the configuration-over-training architecture may then transfer to other computer-science schools.

\section{Conclusion}
\label{sec:conclusion}

Novedu is a design case of how a school can move from exploratory AI use toward teacher-configurable, privacy-aware tutoring. The survey answers RQ1 with four requirements; the architecture answers RQ2 by making alignment a teacher-owned configuration problem rather than a model-training one; and the harness answers RQ3 tentatively: on one topic and a small sample, a small in-region model did not fall behind frontier baselines in holding the Socratic stance, provided teachers supply detailed prompts.

\bibliographystyle{splncs03}
\bibliography{references}

\end{document}
